\documentclass[11pt]{article}
\usepackage[a4paper,margin=1in]{geometry}
\usepackage{amsmath,amssymb,amsfonts}
\usepackage{mathtools}
\usepackage{hyperref}
\usepackage{enumitem}
\hypersetup{colorlinks=true,linkcolor=blue,citecolor=blue,urlcolor=blue}
\setlist{nosep}

\title{Stochastic Long-Maturity Boundary States and the Completed Boundary Jacobian}
\author{Sachin Oza}
\date{\today}

\begin{document}
\maketitle

\begin{abstract}
This paper extends the IRFR boundary-field framework from the short end of the maturity manifold to the stochastic long end. Earlier work introduced the IRFR field representation and showed that short-end boundary closures generate the linear-exponential repeated-root hump \(\tau e^{-a^2\tau}\). The present paper distinguishes the equilibrium long-run anchor \(x_\infty\) from the realised long-boundary state \(X_t\), derives a minimal multiplicative long-boundary law, and embeds the long boundary in a completed three-loading maturity representation. In the canonical bridge case, the field is generated by the Jacobian \(J(\tau)=(e^{-a^2\tau},\tau e^{-a^2\tau},1-e^{-a^2\tau})\). The first loading transports the short boundary, the second is the repeated-root belly mode inherited from the short-boundary construction, and the third transports the long boundary. The paper separates this invariant maturity geometry from the stochastic rank of the state vector. It then records a natural two-direction stochastic specification, with correlated short/belly and long-boundary Brownian directions, as the simplest dynamic layer consistent with the completed Jacobian. The empirical validation of that rank structure is left to subsequent work.
\end{abstract}


\section{Introduction}

The invariant risk-free rate (IRFR) boundary-field framework begins from a simple premise: the term structure is not merely a collection of independent maturities, but a maturity field constrained by boundary conditions. The first paper in this sequence developed the basic IRFR field representation and the associated maturity transport geometry. The second paper studied the short end of the maturity manifold, showing that Neumann/free, Dirichlet, and Robin/mixed boundary closures generate distinct admissible volatility structures. In particular, the repeated-root loading
\[
\tau e^{-a^2\tau}
\]
emerges as the boundary-induced hump produced when stochastic exposure is displaced away from a fixed or partially anchored short endpoint.

This paper completes the corresponding long-end side of the construction. The earlier IRFR representation treats the asymptotic long-maturity anchor as fixed. That is appropriate as an equilibrium reference, but too restrictive if the realised long end of the curve can move persistently. The present paper therefore distinguishes the equilibrium anchor \(x_\infty\) from the realised stochastic long-boundary state \(X_t\). This separation permits persistent long-end displacement while preserving the boundary-field logic of the earlier papers.

The central result is a completed three-loading maturity Jacobian. In the canonical bridge case, the rate field can be written
\[
x_t(\tau)
=
e^{-a^2\tau}x_{t,t}
+
\tau e^{-a^2\tau}B_t
+
\left(1-e^{-a^2\tau}\right)X_t.
\]
The first loading transports the short boundary into maturity space. The second is the repeated-root belly mode already derived in the short-boundary paper. The third transports the stochastic long boundary into maturity space. The completed Jacobian is therefore
\[
J(\tau)
=
\left(
 e^{-a^2\tau},
 \tau e^{-a^2\tau},
 1-e^{-a^2\tau}
\right).
\]

This representation is deliberately separated from any particular stochastic rank assumption. The minimal multiplicative long-boundary law developed below gives a one-direction diffusion for \(X_t\). Once the belly amplitude is promoted to a calendar-time state, however, its dynamics belong to the chosen empirical or pricing layer. The invariant theoretical object supplied by this paper is the maturity geometry; the stochastic rank required to describe calendar-time variation is an empirical and modelling question. For that reason, the paper also records a natural two-direction implementation in which the short/belly component and the long-boundary component are driven by correlated Brownian shocks. This implementation clarifies the natural stochastic reading of the completed Jacobian without making empirical rank a premise of the theory.

The relation to earlier empirical volatility-hump models is therefore constructive. Linear-exponential volatility specifications, such as those associated with Ritchken and Chuang, may be read retrospectively as identifying the short-side part of the boundary geometry: the decaying short loading and the repeated-root hump. The present framework preserves that empirically useful geometry, gives it a boundary-condition interpretation through the short-boundary paper, and completes it by adding the stochastic long-boundary loading.

The first paper in this sequence introduced the IRFR boundary-field representation. The second paper classified admissible short-end boundary behaviour and showed how fixed or partially fixed short-rate boundaries generate the linear-exponential volatility hump. The present paper extends that framework to the opposite end of the maturity manifold. Neumann/free, Dirichlet, and Robin/mixed boundary regimes determine how the maturity field behaves at the instantaneous short end. The terminology ``free'' refers to the local boundary condition or derivative constraint. It does not mean that the realised short endpoint is absent. The instantaneous short endpoint remains a boundary state of the field. In the completed representation, the short end is pinned to the realised short boundary, the long end is pinned to the realised long boundary, and the interior of the maturity curve is free to deform between them.

This paper completes the boundary-field picture by relaxing the opposite end of the maturity manifold. Instead of forcing the long-maturity boundary to coincide pathwise with the fixed equilibrium anchor \(x_\infty\), we distinguish the equilibrium anchor from the realised asymptotic boundary state. The central distinction is
\[
x_\infty
\quad\text{as the equilibrium long-run anchor,}
\]
and
\[
X_t
\quad\text{as the realised long-maturity boundary state.}
\]
The boundary displacement is
\[
Y_t := X_t-x_\infty.
\]

The construction may be read in two layers. In the minimal physical subclass developed first, the displacement \(Y_t\) is carried by the same admissible IRFR risk direction as the finite-maturity field. This preserves the one-risk-direction admissibility structure. In the completed boundary-field representation, however, the essential object is the maturity-space Jacobian generated by the short boundary, the interior deformation, and the long boundary. This Jacobian is the object that subsequent empirical work can test without re-estimating the theory's geometry, while the stochastic rank of the state vector is left to the chosen dynamic or empirical layer.

\section{Motivation}

In the fixed long-boundary formulation, the maturity field is anchored asymptotically at \(x_\infty\). Finite-maturity deviations may be transported along the curve, but the limiting long end remains fixed:
\[
\lim_{\tau\to\infty} x_{t,T}=x_\infty,
\qquad \tau=T-t.
\]
This is appropriate when the long-run equilibrium anchor is treated as both the economic regime reference and the realised long-end state. It is too restrictive, however, if persistent long-run level shifts or long-end volatility are to be represented within the same boundary-field architecture.

The stochastic long-boundary extension separates these roles. The parameter \(x_\infty\) remains the equilibrium anchor. The process \(X_t\) is the realised long-end boundary state. When \(X_t\neq x_\infty\), the asymptotic boundary of the field is displaced from equilibrium. Even before formal empirical testing, a cursory observation of yield-curve behaviour suggests that the realised far end of the curve is not pathwise pinned to a deterministic constant. Long maturities exhibit persistent level movement and low-frequency regime shifts. Treating \(x_\infty\) as an equilibrium reference while allowing \(X_t\) to be stochastic is therefore a natural relaxation of the fixed-boundary representation.

The purpose of this paper is therefore to show that stochastic long-end boundary motion can be incorporated without abandoning the boundary-field architecture. In the minimal multiplicative subclass, the long-boundary displacement can be carried by a single admissible risk direction. In the completed boundary-field representation, however, the long-boundary state may also be assigned its own Brownian direction, correlated with the short/belly direction. The short-end boundary classification remains intact; the long-boundary layer supplies the missing asymptotic state coordinate.

A useful way to state the completed boundary picture is therefore:
\[
\text{short boundary pinned,}
\qquad
\text{long boundary pinned,}
\qquad
\text{interior deformation free.}
\]
The short boundary supplies the realised instantaneous endpoint. The long boundary supplies the realised asymptotic endpoint. The nontrivial curve-shape degree of freedom is the interior, or belly, mode that vanishes at both endpoints.

\section{Boundary decomposition}

Let
\[
\tau=T-t.
\]
For a given short-end boundary branch \(j\), define the branch transport reference
\[
b_j(\tau)
=
x_\infty\left(1-e^{-r_j a^2\tau}\right).
\]
Here \(r_j\) is the branch parameter associated with the short-end boundary regime. The stochastic long-boundary loading is defined by
\[
L(\tau)
=
1-e^{-a^2\tau}.
\]
This loading satisfies
\[
L(0)=0,
\qquad
\lim_{\tau\to\infty}L(\tau)=1.
\]

The stochastic long-boundary decomposition is
\[
x^j_{t,T}
=
b_j(\tau)
+
L(\tau)(X_t-x_\infty)
+
Z^j_{t,T}.
\]
Equivalently,
\[
Z^j_{t,T}
:=
x^j_{t,T}
-
b_j(\tau)
-
L(\tau)(X_t-x_\infty).
\]
The residual \(Z^j_{t,T}\) is the finite-maturity curve-shape component after removing both the branch transport reference and the transported long-boundary displacement.

At the short end,
\[
b_j(0)=0,
\qquad
L(0)=0,
\]
so
\[
x^j_{t,t}=Z^j_{t,t}.
\]
At the long end, provided \(Z^j_{t,T}\to0\) as \(\tau\to\infty\),
\[
\lim_{\tau\to\infty}x^j_{t,T}
=
x_\infty+(X_t-x_\infty)
=
X_t.
\]
Thus the decomposition preserves the short-end state while replacing the fixed long-end limit \(x_\infty\) with the realised stochastic boundary state \(X_t\).

\section{Single-risk-direction admissibility}

The stochastic boundary branch is defined by requiring the full displacement from the branch transport reference to be carried by the unique admissible IRFR risk direction:
\[
x^j_{t,T}-b_j(\tau)
=
L(\tau)(X_t-x_\infty)+Z^j_{t,T}.
\]
The boundary displacement and the finite-maturity residual therefore do not generate separate Brownian directions. They enter the same drift--diffusion loading.

Define
\[
f_j(t,T)
=
a^2\left[x^j_{t,T}-b_j(\tau)\right],
\]
and
\[
g_j(t,T)
=
a\left[x^j_{t,T}-b_j(\tau)\right].
\]
Then, whenever the effective loading is non-zero,
\[
\frac{f_j(t,T)}{g_j(t,T)}
=
a.
\]
The physical market signal-to-noise ratio is therefore preserved:
\[
\mathrm{MSNR}^{\mathbb P}_j
=
\frac{f_j(t,T)}{\sqrt{2}\,g_j(t,T)}
=
\frac{a}{\sqrt{2}}.
\]

Under the IRFR field convention,
\[
dx^j_{t,T}
=
-f_j(t,T)\,dt
+
\sqrt{2}\,g_j(t,T)\,dW_t,
\]
the stochastic long-boundary branch becomes
\[
dx^j_{t,T}
=
-a^2\left[x^j_{t,T}-b_j(\tau)\right]dt
+
\sqrt{2}a\left[x^j_{t,T}-b_j(\tau)\right]dW_t.
\]
The stochastic boundary changes the level of the effective loading, but it does not change the admissible risk direction.

\section{Asymptotic boundary dynamics}

Taking the long-maturity limit of the field equation gives the boundary-mode dynamics. Since
\[
x^j_{t,T}\to X_t,
\qquad
b_j(\tau)\to x_\infty
\quad
\text{as }\tau\to\infty,
\]
the limiting equation is
\[
dX_t
=
-a^2(X_t-x_\infty)\,dt
+
\sqrt{2}a(X_t-x_\infty)\,dW_t.
\]
Equivalently, for
\[
Y_t:=X_t-x_\infty,
\]
we have
\[
dY_t
=
-a^2Y_t\,dt
+
\sqrt{2}aY_t\,dW_t.
\]

A boundary mode, in the sense used here, is the limiting maturity-field excitation carried by the asymptotic boundary state rather than by an interior curve-shape residual. The long-boundary displacement \(Y_t\) is therefore a multiplicative boundary mode. For \(Y_{t_0}\neq0\),
\[
Y_t
=
Y_{t_0}
\exp\left(
-2a^2(t-t_0)
+
\sqrt{2}a(W_t-W_{t_0})
\right).
\]
The sign of the displacement is preserved. If \(Y_{t_0}=0\), then \(Y_t=0\) thereafter and \(X_t=x_\infty\). Thus, in this multiplicative long-boundary subclass, departures from the equilibrium long-end anchor must either be present as an initial boundary displacement or introduced through a richer additive or regime-switching boundary mechanism.

This is a modelling implication of the multiplicative construction, and it is a substantive restriction rather than a purely technical detail. The stochastic long boundary represents persistence and propagation of an existing long-run displacement. It does not by itself generate a new displacement from the exactly anchored state \(X_t=x_\infty\). In empirical applications this restriction is usually benign when the estimated long boundary is not exactly equal to the equilibrium anchor. Nevertheless, it is a testable assumption of the minimal multiplicative subclass: if the data require repeated departures from an exactly anchored long state, then the boundary law must be enlarged by an additive innovation, a regime-switching component, or a higher-rank long-boundary mechanism. The present paper deliberately retains the minimal one-risk-direction construction because it isolates the geometry of persistent long-end displacement.

\section{Generator corroboration}

The same boundary law is recovered from the limiting infinitesimal generator. In the long-maturity limit, write
\[
y=X-x_\infty.
\]
For smooth test functions \(\varphi\), the physical generator associated with the boundary process is
\[
(\mathcal A\varphi)(X)
=
-a^2(X-x_\infty)\varphi_X(X)
+
a^2(X-x_\infty)^2\varphi_{XX}(X).
\]
Applying \(\mathcal A\) to \(\varphi(X)=X\) gives the drift
\[
-a^2(X-x_\infty).
\]
The second-order term identifies the diffusion variance. For \(X_0=X\),
\[
\mathcal A\left[(X-X_0)^2\right]\bigg|_{X=X_0}
=
2a^2(X_0-x_\infty)^2.
\]
Thus
\[
\sigma^2(X)
=
2a^2(X-x_\infty)^2,
\]
with signed one-risk-direction choice
\[
\sigma(X)
=
\sqrt{2}a(X-x_\infty).
\]
The generator calculation therefore corroborates the asymptotic boundary SDE without introducing a second Brownian source or changing the stochastic basis.

\section{Expected boundary reversion}

For fixed maturity \(T\), let
\[
m_j(t,T):=\mathbb E[x^j_{t,T}].
\]
The conditional mean satisfies
\[
\frac{d m_j(t,T)}{dt}
=
-a^2\left[m_j(t,T)-b_j(T-t)\right].
\]
The equilibrium mean curve in maturity coordinates is
\[
x^j_e(\tau)
=
x_\infty
-
\frac{x_\infty}{1+r_j}e^{-r_j a^2\tau}.
\]
Thus,
\[
\mathbb E[x^j_{t,T}]
=
x^j_e(T-t)
+
e^{-a^2(t-t_0)}
\left[
x^j_{t_0,T}
-
x^j_e(T-t_0)
\right].
\]
Taking the long-maturity limit gives
\[
\lim_{T-t\to\infty}\mathbb E[x^j_{t,T}]
=
x_\infty
+
(X_{t_0}-x_\infty)e^{-a^2(t-t_0)}.
\]
The expected realised long boundary therefore mean-reverts toward the equilibrium anchor \(x_\infty\) at the physical adjustment speed \(a^2\).

The branch transport reference \(b_j(\tau)\) and the equilibrium mean curve \(x^j_e(\tau)\) should not be confused. The former is the instantaneous reference appearing inside the local drift loading. The latter is the stationary first-moment profile after accounting for maturity ageing, \(\tau=T-t\). Except in the degenerate case \(r_j=0\), these two objects differ.

\section{Long-end variance floor}

The instantaneous variance rate of the stochastic long-boundary field is
\[
v_j(t,T)
:=
\frac{d}{dt}\operatorname{Var}_t(x^j_{t,T})
=
2a^2\left[x^j_{t,T}-b_j(\tau)\right]^2.
\]
Taking the long-maturity limit yields
\[
\lim_{\tau\to\infty}v_j(t,T)
=
2a^2(X_t-x_\infty)^2.
\]
This is the main observable implication of the stochastic long-boundary extension. In the fixed-boundary model, the long-end loading vanishes and the local variance collapses to zero. In the stochastic-boundary model, long-end variance remains positive whenever the realised long boundary differs from the equilibrium anchor. Conversely, the same multiplicative structure predicts that the long-end variance rate should become small when \(X_t\) lies close to \(x_\infty\). The variance floor is therefore not an unconditional constant floor. It is a displacement-dependent floor generated by the stochastic long boundary.

This gives the chapter its most direct empirical implication: the long-end variance should be organised by the squared boundary displacement \((X_t-x_\infty)^2\), scaled by \(2a^2\). If observed long-maturity variance behaves instead like a level-independent constant, then the minimal multiplicative boundary law is incomplete and must be supplemented by an additional long-end noise source. The formulation below therefore supplies a sharp diagnostic rather than a purely descriptive volatility term.

The long-end variance floor is therefore tied to three interpretable quantities:
\[
a^2,
\qquad
x_\infty,
\qquad
X_t.
\]
The adjustment speed \(a^2\) controls propagation and reversion. The anchor \(x_\infty\) represents the equilibrium long-run regime. The realised boundary \(X_t\) represents the market-implied long-run state. Persistent long-end volatility is therefore represented by the distance between the realised long-run boundary and the equilibrium regime anchor, rather than by an additional curve-fitting factor.


\section{Curve representation and completed geometry}

The stochastic boundary decomposition becomes testable once a finite-maturity residual loading is specified. The completed boundary-field representation uses the endpoint-preserving interior loading
\[
Z^j_{t,T}
=
x^j_{t,t}e^{-a^2\tau}
+
B^j_t\tau e^{-a^2\tau}.
\]
The first term transports the realised short boundary into maturity space. The second term is the interior, or belly, deformation. Its functional form is not introduced here as a new curve-fitting device. It is the same repeated-root loading derived in the short-boundary paper, where a fixed or partially anchored short boundary displaces stochastic exposure away from the endpoint and into the interior of the curve. In the present long-boundary chapter the same mode reappears as the finite-maturity bridge between a pinned short endpoint and a stochastic long endpoint.

The coefficient \(B^j_t\) should therefore be read as the amplitude of the boundary-induced belly mode. In a pure short-end Dirichlet or Robin closure this amplitude is tied to the corresponding short-boundary condition and moment-transport calculation. In the completed bridge representation it may also be treated as a calendar-time state coefficient when the empirical or pricing layer promotes the belly amplitude to a stochastic state. The loading itself is fixed by the boundary geometry.

The belly loading satisfies
\[
\tau e^{-a^2\tau}=0
\quad\text{at }\tau=0,
\qquad
\tau e^{-a^2\tau}\to0
\quad\text{as }\tau\to\infty,
\]
and attains its maximum at
\[
\tau=\frac{1}{a^2}.
\]
Thus it is the boundary-induced free mode between the two pinned endpoints.

Substituting into the stochastic boundary decomposition gives
\[
x^j_{t,T}
=
x_\infty\left(1-e^{-r_j a^2\tau}\right)
+
\left(1-e^{-a^2\tau}\right)(X_t-x_\infty)
+
x^j_{t,t}e^{-a^2\tau}
+
B^j_t\tau e^{-a^2\tau}.
\]
Equivalently,
\[
x^j_{t,T}
=
x^j_{t,t}e^{-a^2\tau}
+
B^j_t\tau e^{-a^2\tau}
+
X_t\left(1-e^{-a^2\tau}\right)
+
x_\infty\left(e^{-a^2\tau}-e^{-r_j a^2\tau}\right).
\]

The branch \(r_j=1\) is the canonical bridge case. In this case the explicit equilibrium-anchor loading cancels from the static cross-sectional representation, leaving
\[
x_{t,T}
=
x_{t,t}e^{-a^2\tau}
+
B_t\tau e^{-a^2\tau}
+
X_t\left(1-e^{-a^2\tau}\right).
\]
This is the completed two-boundary curve decomposition. Its three loadings are
\[
\underbrace{e^{-a^2\tau}}_{\text{short-boundary loading}},
\qquad
\underbrace{\tau e^{-a^2\tau}}_{\text{interior/free loading}},
\qquad
\underbrace{1-e^{-a^2\tau}}_{\text{long-boundary loading}}.
\]

The short-boundary loading is one at the short end and decays away from it. The long-boundary loading is zero at the short end and converges to one at long maturity. The interior loading vanishes at both boundaries and governs the belly of the curve. This is the long-boundary continuation of the short-boundary repeated-root mechanism: the same \(\tau e^{-a^2\tau}\) hump generated by the short-end closure becomes the endpoint-preserving interior deformation of the two-boundary field. The theory is therefore no longer merely an exponential bridge between two endpoints. It is a two-boundary field with a boundary-induced interior deformation.

The equilibrium anchor \(x_\infty\) still enters through the boundary dynamics, expected reversion, and long-end variance floor. In the canonical bridge representation it no longer appears as a separate visible loading in the static curve. Calendar-time variation is carried by the state coefficients \(x_{t,t}\), \(B_t\), and \(X_t\), while maturity-space geometry is carried by the fixed loading triple above.

\section{Branch specialisations}

The stochastic long-boundary machinery is common to all short-end boundary branches. The branch label \(j\) changes the short-end admissibility condition and the parameter \(r_j\) appearing in the transport reference \(b_j(\tau)\). It does not change the multiplicative long-boundary law, the MSNR preservation result, or the long-end variance floor derived above.

For the Neumann/free branch,
\[
r_j=r.
\]
For the Dirichlet branch,
\[
r_D=\frac{\lambda_D^2}{a^2}.
\]
For the Robin/mixed branch,
\[
r_R=\frac{\lambda_R^2}{a^2}.
\]
Across all three branches,
\[
b_j(\tau)
=
x_\infty\left(1-e^{-r_j a^2\tau}\right).
\]
The general curve representation has already been written above as
\[
x^j_{t,T}
=
x^j_{t,t}e^{-a^2\tau}
+
B^j_t\tau e^{-a^2\tau}
+
X_t\left(1-e^{-a^2\tau}\right)
+
x_\infty\left(e^{-a^2\tau}-e^{-r_j a^2\tau}\right).
\]
The final term vanishes only in the canonical bridge case \(r_j=1\). For \(r_j\neq1\), the equilibrium anchor remains visible in the static curve as a branch-specific displacement between the long-boundary loading and the branch transport loading. Thus the cancellation of \(x_\infty\) in the canonical bridge case is not a hidden assumption of the general theory; it is the special simplification that makes the \(r_j=1\) branch the cleanest object for empirical testing.

The Dirichlet and Robin amplitudes are not inherited from the Neumann/free branch. They are fixed by the branch-specific equilibrium or moment-transport calculation developed in the short-boundary papers. In the present paper, the long-boundary mechanism is layered on top of those branch-specific short-end closures.


\section{Volatility geometry and the boundary Jacobian}

The completed boundary-field representation separates maturity geometry from state dynamics. In the canonical bridge case, define the maturity-space Jacobian
\[
J(\tau)
=
\left(
 e^{-a^2\tau},
 \tau e^{-a^2\tau},
 1-e^{-a^2\tau}
\right).
\]
For the state vector
\[
S_t=
\left(x_{t,t},B_t,X_t\right)^{\top},
\]
the curve may be written compactly as
\[
x_t(\tau)=J(\tau)S_t.
\]
The three entries of \(J(\tau)\) have distinct boundary meanings. The first transmits the short boundary into maturity space. The second is the repeated-root belly mode inherited from the short-end boundary closure and now serving as the endpoint-preserving interior deformation. The third transmits the stochastic long boundary into maturity space.

The three loadings are not asserted to be orthogonal. They are structural boundary coordinates rather than statistical principal components. What is required is linear independence. Indeed, if
\[
c_1e^{-a^2\tau}+c_2\tau e^{-a^2\tau}+c_3(1-e^{-a^2\tau})=0
\]
for every \(\tau\geq0\), then letting \(\tau\to\infty\) gives \(c_3=0\). The remaining identity becomes
\[
(c_1+c_2\tau)e^{-a^2\tau}=0,
\]
which forces \(c_1=c_2=0\). Thus the Jacobian loadings form a linearly independent boundary-generated coordinate system. Orthogonal statistical factors may be obtained by diagonalising or rotating covariance operators within this span, but the structural content lies in the boundary loadings themselves.

This representation should not be read as imposing a particular stochastic rank on the state vector. The multiplicative long-boundary construction above gives a minimal one-direction diffusion for \(X_t\). The belly amplitude \(B_t\) has already appeared in the short-boundary theory as a boundary-induced coefficient. When it is promoted from a branch-specific amplitude to an empirical or pricing state, its calendar-time dynamics must be specified by that chosen dynamic layer. In general, write
\[
dS_t=\mu_t\,dt+Q_t\,dW_t,
\]
where \(W_t\) may be one-dimensional or multidimensional. The instantaneous covariance matrix of the state vector is then
\[
\Sigma_t=Q_tQ_t^{\top},
\]
and the corresponding maturity covariance kernel is
\[
K_t(\tau,\sigma)
=
J(\tau)\Sigma_tJ(\sigma)^{\top}.
\]

Thus the invariant object supplied by the boundary theory is the maturity-space Jacobian \(J(\tau)\), not a fixed stochastic rank. A one-direction, two-direction, or higher-rank volatility structure corresponds to different restrictions on \(\Sigma_t\). What remains invariant is the maturity geometry: the short-boundary direction, the repeated-root belly direction, and the long-boundary direction. This is the final theoretical object supplied by the stochastic long-boundary construction and the natural object to be tested in the empirical paper.

\subsection{A natural two-direction stochastic specification}

Although the boundary Jacobian does not itself impose a stochastic rank, it suggests a natural low-rank dynamic implementation. The short-boundary loading and the repeated-root belly mode are inherited from the short-end boundary closure. It is therefore natural to group them into a short/belly shock. The long-boundary loading represents the realised asymptotic state and may be driven by a distinct long-boundary shock. Allow these shocks to be correlated:
\[
dW^S_t\,dW^X_t=\rho_t\,dt.
\]
A simple two-direction volatility layer is then
\[
dx_t(\tau)
=
g_S(t,\tau)\,dW^S_t
+
g_X(t,\tau)\,dW^X_t,
\]
with
\[
g_S(t,\tau)
=
\nu^r_t e^{-a^2\tau}
+
\nu^B_t\tau e^{-a^2\tau},
\qquad
g_X(t,\tau)
=
\nu^X_t\left(1-e^{-a^2\tau}\right).
\]
Here the coefficients \(\nu^r_t\), \(\nu^B_t\), and \(\nu^X_t\) may be deterministic, state-dependent, or estimated from data, depending on the empirical or pricing layer. This specification is not required by the boundary geometry; it is the simplest stochastic implementation consistent with the completed Jacobian and with the interpretation that the belly mode is short-boundary induced while the long-boundary state has its own persistent source of risk.

The corresponding covariance kernel is
\[
\begin{aligned}
K_t(\tau,\sigma)
&=
g_S(t,\tau)g_S(t,\sigma)
+
g_X(t,\tau)g_X(t,\sigma) \\
&\quad+
\rho_t\{g_S(t,\tau)g_X(t,\sigma)+g_X(t,\tau)g_S(t,\sigma)\}.
\end{aligned}
\]
This formula records the natural two-Brownian reading of the completed field: one direction for the short/belly boundary geometry and one direction for the stochastic long boundary, with correlation \(\rho_t\). The purpose of the present paper is to define the geometry and the admissible dynamic form. Whether the data require one, two, or more stochastic directions is an empirical question taken up in subsequent work.

\section{Relation to risk-neutral pricing}

The stochastic long-boundary extension is a physical-field construction. It identifies how a realised long-maturity boundary state can move within the same admissible IRFR risk direction. The following chapter introduces the risk-neutral pricing layer.

This ordering is important. The money-market account is generated by the realised short-end roll \(x_{t,t}\). The stochastic long boundary \(X_t\) does not directly fund the bank account. Instead, it enters bond prices through the IRFR strip
\[
\{x_{t,u}:u\in[t,T]\}.
\]
When the long boundary is stochastic, the bond-price functional inherits dependence on the realised asymptotic state. A full no-arbitrage treatment must therefore trace \(X_t\) through the risk-neutral drift restriction and the rolling-maturity transport equation.

The present paper does not impose that pricing restriction. It completes the physical boundary-field architecture. The next chapter adds the money-market numeraire, the martingale condition for discounted bond prices, and the change from the physical measure to the pricing measure.


\section{Summary}

The fixed-boundary IRFR model identifies the long end with the equilibrium anchor \(x_\infty\). The stochastic long-boundary extension separates this anchor from the realised long-end state \(X_t\). The resulting displacement
\[
Y_t=X_t-x_\infty
\]
is carried by the same admissible IRFR risk direction as the finite-maturity field in the minimal physical subclass. The completed field, however, does not require the belly amplitude and the long boundary to share a single Brownian source.

The construction preserves the physical loading relation
\[
\frac{f_j}{g_j}=a,
\]
and hence preserves the one-risk-direction MSNR structure. The asymptotic boundary dynamics are
\[
dX_t
=
-a^2(X_t-x_\infty)\,dt
+
\sqrt{2}a(X_t-x_\infty)\,dW_t.
\]
The long-end variance floor is
\[
\lim_{\tau\to\infty}v_j(t,T)
=
2a^2(X_t-x_\infty)^2.
\]

The completed boundary-field geometry is clearest in the canonical bridge branch \(r_j=1\), where the curve is
\[
x_t(\tau)
=
x_{t,t}e^{-a^2\tau}
+
B_t\tau e^{-a^2\tau}
+
X_t\left(1-e^{-a^2\tau}\right).
\]
This representation has a precise boundary interpretation. The short endpoint is pinned to the realised short boundary. The long endpoint is pinned to the realised long boundary. The interior is free to deform through the repeated-root belly mode \(\tau e^{-a^2\tau}\), the same boundary-induced hump derived in the short-boundary paper, which vanishes at both endpoints.

Thus the boundary-field architecture has two complementary layers. Short-end boundary regimes determine admissibility at \(\tau=0\). The stochastic long-boundary state determines persistent asymptotic displacement as \(\tau\to\infty\). The interior loading supplies the finite-maturity deformation between those pinned endpoints. Together they produce the maturity-space Jacobian
\[
J(\tau)
=
\left(
 e^{-a^2\tau},
 \tau e^{-a^2\tau},
 1-e^{-a^2\tau}
\right),
\]
which governs both the level geometry and the volatility geometry of the field. Calendar-time variation moves the state coefficients; maturity-time reveals the invariant boundary geometry. This is the completed theoretical structure to be carried into empirical testing.


\begin{thebibliography}{9}

\bibitem{OzaIRFR}
Oza, S. (2025a).
\emph{Invariant Risk-Free Rate Fields and the Boundary Representation}.
Working paper.

\bibitem{OzaShortBoundary}
Oza, S. (2025b).
\emph{Short-End Boundary Conditions and the Volatility Hump in Interest-Rate Fields}.
Working paper.

\bibitem{RitchkenChuang}
Ritchken, P., and Chuang, I. (1999).
Interest rate option pricing with volatility humps.
Working paper / published reference to be inserted.

\end{thebibliography}
\end{document}
